\documentclass[10pt,twocolumn]{article}
\usepackage[margin=1.8cm]{geometry}
\usepackage{amsmath}
\usepackage{lmodern}
\usepackage[T1]{fontenc}
\title{A Low-Power Bandgap Reference with Curvature Correction}
\author{H. Trinh \and A. Example}
\date{November 2026}
\begin{document}
\maketitle
\begin{abstract}
We present a sub-microwatt bandgap voltage reference for implantable neural interfaces.
The circuit achieves a temperature coefficient of 4.2 ppm/K from $-20$ to $85\,^\circ$C
while drawing only 380 nA from a 1.2 V supply.
\end{abstract}
\section{Introduction}
Voltage references are a fundamental building block of every mixed-signal system.
In neural recording front ends, the reference sets the full-scale range of the
analog-to-digital converter and therefore directly limits the achievable
resolution. Power is the dominant constraint, because the heat dissipated by
an implant must stay below the limit set by tissue safety.
Classic bandgap designs combine a voltage proportional to absolute temperature
(PTAT) with the base-emitter voltage of a bipolar transistor, which is
complementary to absolute temperature (CTAT).
\section{Circuit Design}
The output voltage of a first-order bandgap is
\begin{equation}
V_{\mathrm{REF}} = V_{BE} + \frac{R_2}{R_1}\, V_T \ln N ,
\end{equation}
where $V_T = kT/q$ is the thermal voltage and $N$ is the emitter area ratio.
The residual curvature of $V_{BE}(T)$ limits the temperature coefficient to
roughly 20 ppm/K. We cancel this curvature with a small current that is
proportional to $T \ln T$, generated by a second bipolar pair biased with a
temperature-independent current.
\subsection{Operational Amplifier}
A two-stage amplifier with Miller compensation forces the two branch voltages to
be equal. Its offset is the main source of spread, so we use chopping at 2 kHz.
\section{Measurement Results}
Twenty samples were measured over temperature. Table~\ref{tab:res} summarises
the results and compares them with prior work.
\begin{table}[h]
\centering
\begin{tabular}{lccc}
\hline
 & This work & [1] & [2] \\
\hline
Supply (V) & 1.2 & 1.8 & 1.0 \\
Power (nW) & 456 & 1200 & 520 \\
TC (ppm/K) & 4.2 & 12 & 9.5 \\
Area (mm$^2$) & 0.031 & 0.08 & 0.045 \\
\hline
\end{tabular}
\caption{Comparison with prior work.}
\label{tab:res}
\end{table}
\section{Conclusion}
The proposed reference combines low power with low temperature drift, which
makes it well suited to battery-free implants. Future work will address
trimming and the supply rejection at high frequencies.
\begin{thebibliography}{9}
\bibitem{a} A. Author, ``A nanowatt reference,'' \emph{JSSC}, 2019.
\bibitem{b} B. Author, ``Curvature compensation revisited,'' \emph{TCAS-I}, 2022.
\end{thebibliography}
\end{document}
